%%=============================================================================
%% Projectgegevens 
%%=============================================================================

%% ---------- Jij ------------------------------------------------------------
\newcommand{\studentnaam}{Sarah Turner Burkitt}
\newcommand{\studentemail}{sarah.turnerburkitt@student.hogent.be}

%% Wordt gebruikt als bestandsnaam van je PDF's, bijvoorbeeld
%% Familienaam-Voornaam-rapport.pdf. Geen spaties of leestekens gebruiken.
\newcommand{\bestandsnaam}{TurnerBurkitt-Sarah}

%% ---------- Je project -----------------------------------------------------
\newcommand{\projecttitel}{Gerichte configuratie van automatisch gegenereerde output}
\newcommand{\projectondertitel}{Uitbreiding van document- en consultatietemplates binnen Vesalius.ai}
%% Een of twee zinnen. Wordt gebruikt op de poster en in het voorstel.
\newcommand{\projectpitch}{%
    Dit graduaatsproject breidt het Vesalius-platform uit met gerichtere configuratie van automatisch gegenereerde output: 
    documenttemplates kunnen aan specifieke artsen worden gekoppeld en Scribe-consultatemplates kunnen worden afgestemd op patiënt- 
    of artsgerichte output.%
}

%% ---------- Waar je het doet ------------------------------------------------
\newcommand{\bedrijfsnaam}{Vesalius.ai}

%% Je mentor op de werkvloer (het bedrijf)
\newcommand{\mentornaam}{Bart D'Hoore}
\newcommand{\mentoremail}{Bart.D'Hoore@endare.be}

%% Je begeleider van de opleiding (HOGENT)
\newcommand{\begeleidernaam}{Luc Vervoort}
\newcommand{\begeleideremail}{luc.vervoort@hogent.be}

%% ---------- Je repository ---------------------------------------------------
\newcommand{\repolink}{https://github.com/STB95/Graduaatsproject_TurnerBurkittSarah_26-27}
%% op je poster: hij is een beoordeeld onderdeel van je oplevering.
\newcommand{\repolink}{https://les.lucvervoort.org/sarahturnerburkitt/graduaatsproject}

%% ---------- Vast voor dit opleidingsonderdeel -------------------------------
%% Deze twee hoef je normaal niet aan te passen.
\newcommand{\opleidingnaam}{Graduaat in het Programmeren}
\newcommand{\academiejaarwaarde}{2026-2027}

%% Zoekwoorden voor het voorstel en de poster (3 tot 5)
\newcommand{\projectkeywords}{Vesalius.ai, documenttemplates, Scribe, automatische output, Angular/Laravel, WYSIWYG e-maileditor}
